% TOP_PROBLEM: {"id": 2303004, "title": "Research Problems in Function Theory — Problem 3.4", "category_id": 8, "category": {"id": 8, "name": "analysis", "display_name": "Analysis", "description": "Limits, continuity, calculus, and function theory.", "slug": "analysis", "order_index": 8, "created_at": "2026-07-31T15:26:25.671Z"}, "status": "open"}
% FIRST_POSTED: null
% ENTRY_KIND: statement_only
% Imported from: batch12.tex; UnsolvedMath ID 2303004
% Compile twice: lualatex 2303004.tex
\documentclass[11pt]{article}
\usepackage{fontspec}
\setmainfont{Latin Modern Roman}
\usepackage{amsmath,amssymb,mathtools,unicode-math}
\setmathfont{Latin Modern Math}
\usepackage{xurl,hyperref}
\hypersetup{colorlinks=true,urlcolor=blue,pdftitle={UnsolvedMath Problem 2303004}}
\newcommand{\sref}[2]{\hyperlink{source:#1:#2}{[S#2]}}
\newcommand{\cataloguescope}[1]{\paragraph{Mathematical question.}#1}
\begin{document}
% UnsolvedMath ID 2303004; source catalogue code AMR-022-3004
\setcounter{section}{2303003}
\section{Research Problems in Function Theory — Problem 3.4}\label{2303004}
{\small\sffamily UnsolvedMath ID 2303004 · Analysis}
\subsection{Definitions and mathematical statement}

\paragraph{Shared notation.}$\mathbb N=\{1,2,\ldots\}$, and $\mathbb Z,\mathbb Q,\mathbb R,\mathbb C$ denote the integers, rationals, reals and complex numbers. Any other conventions are specified locally.


A subharmonic function is an upper semicontinuous function with values in $[-\infty,\infty)$, not identically $-\infty$ on a connected component, that satisfies the mean-value inequality on every closed ball in its domain. A real harmonic function is a twice continuously differentiable function $u$ satisfying $\Delta u=0$. Let $\mathbb D=\{|z|<1\}$ and let $E\subset[0,1]$ be a finite union of intervals. Write $A(r,u)=\inf_{|z|=r}u(z)$. Let $u_E$ be the Carleman--Milloux slit function: it is harmonic on $\mathbb D\setminus E$, has boundary values $0$ on $\partial\mathbb D\setminus E$ and $-1$ on $E\cap\mathbb D$, and is extended by $-1$ on the slit. For $0<r<1$ and $0<K<1$ set
\[C_u(r,K)=\{\theta\in[0,2\pi):u(re^{i\theta})<-K\}.\]
Length of this angular set means one-dimensional Lebesgue measure.
\paragraph{Statement recorded in the supplied source.}
Among subharmonic $u\leq0$ on $\mathbb D$ with $A(r,u)\leq-1$ for every $r\in E\cap[0,1)$ (with $A(0,u)=u(0)$), is $|C_u(r,K)|$ minimized only when $u=u_E$? In particular, answer the question for the full slit $E=[0,1]$. \sref{2303004}{2}
\subsection{Short English statement}
Is the Carleman--Milloux slit function the unique minimizer of the angular region below a prescribed negative level, especially for a full radial slit?
\subsection{Sources}
\begin{description}
\item[\hypertarget{source:2303004:1}{S1}] ULAM.AI and the UnsolvedMath Contributors, UnsolvedMath: A Curated Collection of Open Mathematics Problems, record 2303004 (AMR-022-3004). CC BY 4.0. \url{https://huggingface.co/datasets/ulamai/UnsolvedMath}.
\item[\hypertarget{source:2303004:2}{S2}] W. K. Hayman and E. F. Lingham, \emph{Research Problems in Function Theory}, arXiv:1809.07200v2 (2018), Problem 3.4. \url{https://arxiv.org/abs/1809.07200}.
\end{description}
\label{end:2303004}
\end{document}
